\documentclass[aps,prd,twocolumn,superscriptaddress,nofootinbib,10pt]{revtex4-2}

\usepackage{amsmath,amssymb}
\usepackage[colorlinks=true,citecolor=blue,linkcolor=blue,urlcolor=blue]{hyperref}

\begin{document}

\title{Primordial helium bounds on a leptonic electron-mass\\
transition inside the nucleosynthesis window}

\author{Justin Pulford}
\email{pulfordj420@gmail.com}
\affiliation{Unaffiliated}

\date{\today}

\begin{abstract}
We treat the amplitude $\varepsilon$ of a universal fractional shift in the electron mass that
turns on inside the big-bang nucleosynthesis (BBN) temperature window as a free parameter, and ask
what the measured primordial helium mass fraction $Y_p$ implies for it. The shift is switched on
below a critical temperature $T_c$ by a linear ramp,
$\varepsilon_{\rm eff}(T)=\varepsilon\,\max(0,1-T/T_c)$, which deforms the weak rates in the
nuclear network without changing the relativistic degrees of freedom. At the $T_c$ where the
response was measured in a production BBN code ($T_c\simeq 179\,$keV), the helium response is
linear, $dY_p/d\varepsilon = 0.00163$ per percent of $\varepsilon$, and reproduces a direct
windowed run to a few percent of the window itself. Against the Aver \emph{et al.}\ determination
$Y_p=0.2453\pm 0.0034$, that elasticity implies
%
\begin{equation*}
  \varepsilon < 3.2\% \qquad (2\sigma).
\end{equation*}
%
The EMPRESS determination cannot be used for the same purpose: standard BBN with $\varepsilon=0$
already sits $2.9\sigma$ above it, so a measurement in tension with the null hypothesis does not
constrain a deformation of that null. Deuterium is not used for a derivative bound, because its
$\varepsilon$-response is a nonlinear bottleneck quantity. The result is a genuine upper limit on
the amplitude of an electron-mass transition that switches on inside the nucleosynthesis window,
with zero fitted parameters, and no dependence on any unconverged cosmological posterior.
\end{abstract}

\maketitle

% ---------------------------------------------------------------------------
\section{Introduction}\label{sec:intro}

A fractional change in the electron mass at early times is a standard deformation of big-bang
nucleosynthesis and of recombination. In the microwave background it is partly degenerate with
the expansion history and has been studied as a possible contribution to the Hubble
tension~\cite{HartChluba2018,HartChluba2020,SekiguchiTakahashi2021}. At nucleosynthesis the same
deformation acts through the weak rates and the binding-sensitive nuclear network, and the
response matrices for varying constants have long been tabulated~\cite{DentSternWetterich2007}.

Most existing bounds assume a constant shift from some high temperature through the BBN window,
or a power-law variation in redshift. The case of interest here is different: an electron-mass
shift that is \emph{absent} above a critical temperature $T_c$ and turns on below it. That
topology matters because neutron--proton freeze-out and the deuterium bottleneck sit at different
temperatures ($\sim\!800\,$keV and $\sim\!70\,$keV, respectively). A transition that switches on
between them deforms the two epochs unequally.

We do not propose a mechanism for such a transition, and we do not fix $T_c$ or $\varepsilon$ to
any particular model value. The question is purely kinematic: if an electron-mass shift of
amplitude $\varepsilon$ turns on inside the BBN window, what does the measured helium mass fraction
allow? The answer is a single number against one observational determination, and an explicit
statement that the other common determination cannot be used for this purpose.

Section~\ref{sec:setup} fixes the ramp and the free parameters. Section~\ref{sec:helium} reports
the measured helium response. Section~\ref{sec:bound} derives the Aver bound.
Section~\ref{sec:empress} explains why EMPRESS yields no bound.
Section~\ref{sec:notused} states what is deliberately left out. Section~\ref{sec:conclude}
summarises.

% ---------------------------------------------------------------------------
\section{Setup}\label{sec:setup}

Let $m_e^{(0)}$ be the present electron mass. We write a temperature-dependent fractional shift
%
\begin{equation}
  m_e(T) = m_e^{(0)}\,\bigl[1 + \varepsilon_{\rm eff}(T)\bigr],
  \label{eq:me}
\end{equation}
%
with a linear turn-on below a critical temperature $T_c$,
%
\begin{equation}
  \varepsilon_{\rm eff}(T)
  = \varepsilon \,\max\bigl(0,\, 1 - T/T_c\bigr).
  \label{eq:ramp}
\end{equation}
%
Above $T_c$ the shift is off; at $T\to 0$ it saturates at the constant amplitude $\varepsilon$.
Both $\varepsilon$ and $T_c$ are free. The only structural commitment is the ramp shape of
Eq.~\eqref{eq:ramp}, which is the minimal continuous switch-on that vanishes above $T_c$ and is
linear in the order-parameter sense below it. Nothing below depends on why the switch-on has that
form.

The relevant temperature window for a transition that BBN can actually witness is bounded below
by the deuterium bottleneck ($\sim\!70\,$keV) and above by the weak-rate window approaching
neutron--proton freeze-out (of order a few hundred keV). We adopt the interval
%
\begin{equation}
  T_c \in [70,\,500]\,\mathrm{keV}
  \label{eq:window}
\end{equation}
%
as the range over which $T_c$ is regarded as free for the present purpose. Outside that interval
the transition either misses the bottleneck entirely (too low) or reaches freeze-out in full (too
high); both extremes change the interpretation of a helium-only bound and are not the case
studied here.

The shift enters BBN through the electron mass in the weak rates. It does not add relativistic
degrees of freedom: a network run with the ramp at fixed $\varepsilon$ returns $N_{\rm eff}$
unchanged at the level of the code's output~\cite{PRyMordial2023}. The deformation is therefore a
pure rate effect, not an expansion-rate lever.

% ---------------------------------------------------------------------------
\section{Measured helium response}\label{sec:helium}

The bound is built from a production run of the PRyMordial nuclear
network~\cite{PRyMordial2023}, with the ramp of Eq.~\eqref{eq:ramp} applied pointwise to the weak
rates: at each temperature the standard and fully shifted rates are blended by the weight
$w(T)=\max(0,1-T/T_c)$. The run that defines the response was performed at
%
\begin{equation}
  T_c = 179\,\mathrm{keV},\qquad \varepsilon = 1.2543\%,
  \label{eq:run}
\end{equation}
%
on the code's default baryon density. The helium mass fraction moves from the $\varepsilon=0$
baseline to the windowed value as
%
\begin{equation}
  Y_p^{(0)} = 0.246891
  \;\longrightarrow\;
  Y_p^{(\varepsilon)} = 0.248995,
  \label{eq:windowed}
\end{equation}
%
a fractional change of $+0.852\%$.

A four-point scan in $\varepsilon$ at the same $T_c$ establishes that the helium response is
linear to a few percent over the range of interest, with elasticity
%
\begin{equation}
  \frac{d Y_p}{d\varepsilon}
  = 0.00163
  \quad\text{per percent of $\varepsilon$}.
  \label{eq:elast}
\end{equation}
%
Multiplying Eq.~\eqref{eq:elast} by $\varepsilon=1.2543\%$ recovers the window of
Eq.~\eqref{eq:windowed} to $2.8\%$ of the window itself. The linearity is what licenses a
derivative bound; without it the translation from one windowed run into a limit on $\varepsilon$
would not be legitimate.

Helium is almost insensitive to the baryon density, $Y_p \propto \omega_b^{0.04}$. Shifting
$\omega_b$ at the percent level therefore moves the baseline by far less than the observational
error used below, so the bound is baseline-robust against the usual CMB--BBN $\omega_b$
re-inference.

% ---------------------------------------------------------------------------
\section{The Aver bound}\label{sec:bound}

Aver \emph{et al.} report~\cite{Aver2021}
%
\begin{equation}
  Y_p^{\rm Aver} = 0.2453 \pm 0.0034.
  \label{eq:aver}
\end{equation}
%
Writing $Y_p(\varepsilon) = Y_p^{(0)} + 0.00163\,\varepsilon_{\%}$ with $\varepsilon_{\%}$ the
amplitude in percent, and requiring $Y_p(\varepsilon)$ to lie within $n\sigma$ of
Eq.~\eqref{eq:aver}, one finds
%
\begin{align}
  \varepsilon &< 1.11\% \qquad (1\sigma), \label{eq:1sig}\\
  \varepsilon &< 3.20\% \qquad (2\sigma). \label{eq:2sig}
\end{align}
%
The $2\sigma$ statement is the result of this paper:
%
\begin{equation}
  \boxed{\ \varepsilon < 3.2\%\ \ (2\sigma)\ }
  \label{eq:result}
\end{equation}
%
against Aver, at the $T_c$ of Eq.~\eqref{eq:run}, with no fitted parameters.

For orientation: the $\varepsilon=0$ baseline already sits $+0.47\sigma$ above Aver, and the
windowed point of Eq.~\eqref{eq:run} sits at $+1.09\sigma$. The transition is constrained but not
excluded at the amplitude of Eq.~\eqref{eq:run}. The $1\sigma$ column of Eq.~\eqref{eq:1sig} is
reported for completeness; the quotable bound is Eq.~\eqref{eq:result}.

A dense map of the upper edge $\varepsilon_{\rm max}(T_c)$ over the full interval of
Eq.~\eqref{eq:window} is the natural extension of Eq.~\eqref{eq:result}. It requires a grid of
network runs in $T_c$ at fixed observational $Y_p$, which is not reported here. The number in
Eq.~\eqref{eq:result} is the verified bound at the measured $T_c$; it is not a claim that the
same numerical ceiling holds at every point of Eq.~\eqref{eq:window}.

% ---------------------------------------------------------------------------
\section{Why EMPRESS cannot be used}\label{sec:empress}

The EMPRESS collaboration reports a lower helium abundance~\cite{EMPRESS2022},
%
\begin{equation}
  Y_p^{\rm EMPRESS} = 0.2370 \pm 0.0034.
  \label{eq:empress}
\end{equation}
%
Against that determination the same baseline and windowed values read
%
\begin{equation}
  \varepsilon=0:\ +2.91\sigma,
  \qquad
  \varepsilon=1.2543\%:\ +3.53\sigma.
  \label{eq:empress_pull}
\end{equation}
%
Standard BBN with no transition is already $2.9\sigma$ high. The window contributes only
$0.62\sigma$ of the $3.53\sigma$. A measurement that is in tension with the null hypothesis
cannot constrain a deformation of that null: removing the transition entirely leaves $2.9\sigma$
standing. Quoting an EMPRESS-based upper limit on $\varepsilon$ would attribute to the transition
a discrepancy that exists without it.

The observational helium disagreement itself --- Aver and EMPRESS differ by about $1.7\sigma$ on
the same quantity --- is a property of the data, not of the ramp. Until it is resolved, a
helium-only bound on $\varepsilon$ must be reported against Aver and must state why EMPRESS is
not used. That is the entire content of this section.

% ---------------------------------------------------------------------------
\section{What is not used}\label{sec:notused}

\emph{Deuterium.} The deuterium abundance responds to $\varepsilon$, but as a bottleneck residue
rather than an equilibrium abundance. A uniform-$\varepsilon$ derivative under-predicts the
windowed effect, and a four-point scan does not yield a linear elasticity that can be inverted
into a clean upper limit of the kind in Eq.~\eqref{eq:result}. No derivative-based deuterium bound
is claimed. Absolute D/H comparisons that depend on a re-inferred baryon density from a
cosmological fit are outside the scope of this note: they are not needed for
Eq.~\eqref{eq:result}, and they are not used.

\emph{Extra radiation.} An additional relativistic density $\Delta N_{\rm eff}$ moves $Y_p$ and
D/H in the same direction. It is an expansion-rate lever, not a measurement of $\varepsilon$, and
it is not floated here. The bound of Eq.~\eqref{eq:result} is the pure rate response at fixed
standard $N_{\rm eff}$.

\emph{Lithium.} The same windowed run moves $^7$Li/H by $+0.26\%$, a percent-level shift that
neither causes nor cures the cosmological lithium problem. Lithium is recorded as unaffected at
the level relevant to Eq.~\eqref{eq:result} and is not used in the bound.

\emph{Model-fixed $T_c$ or $\varepsilon$.} No central value of $T_c$ or $\varepsilon$ is adopted
as an input. The run of Eq.~\eqref{eq:run} fixes the elasticity; the bound is the upper edge that
Aver allows.

% ---------------------------------------------------------------------------
\section{Conclusion}\label{sec:conclude}

A universal electron-mass shift that turns on inside the BBN window, with free amplitude
$\varepsilon$ and free critical temperature $T_c$ in the nucleosynthesis-relevant range, is
bounded by primordial helium. Against Aver \emph{et al.}, and at the $T_c$ where the helium
elasticity was measured,
%
\begin{equation*}
  \varepsilon < 3.2\% \qquad (2\sigma),
\end{equation*}
%
with zero fitted parameters. The EMPRESS determination cannot be used for this purpose, because
standard BBN is already $2.9\sigma$ high against it before any transition is applied. Deuterium
is excluded from the bound because its $\varepsilon$-response is nonlinear. The result is a
chain-free, single-observable upper limit on the amplitude of a leptonic electron-mass transition
inside the nucleosynthesis window.

\begin{thebibliography}{99}

\bibitem{HartChluba2018}
L.~Hart and J.~Chluba,
New constraints on time-varying fundamental constants from Planck,
Mon.\ Not.\ R.\ Astron.\ Soc.\ \textbf{474}, 1850 (2018);
arXiv:1705.03925.

\bibitem{HartChluba2020}
L.~Hart and J.~Chluba,
Updated fundamental constant constraints from Planck 2018 data and possible relations to the
Hubble tension,
Mon.\ Not.\ R.\ Astron.\ Soc.\ \textbf{493}, 3255 (2020);
arXiv:1912.03986.

\bibitem{SekiguchiTakahashi2021}
T.~Sekiguchi and T.~Takahashi,
Early recombination as a solution to the $H_0$ tension,
Phys.\ Rev.\ D \textbf{103}, 083507 (2021);
arXiv:2007.03381.

\bibitem{DentSternWetterich2007}
T.~Dent, S.~Stern, and C.~Wetterich,
Primordial nucleosynthesis as a probe of fundamental physics parameters,
Phys.\ Rev.\ D \textbf{76}, 063513 (2007);
arXiv:0705.0696.

\bibitem{PRyMordial2023}
A.-K.~Burns, T.~M.~P.~Tait, and M.~Valli,
PRyMordial: the first three minutes, within and beyond the Standard Model,
Eur.\ Phys.\ J.\ C \textbf{84}, 86 (2024);
arXiv:2307.07061.

\bibitem{Aver2021}
E.~Aver, D.~A.~Berg, K.~A.~Olive, R.~W.~Pogge, J.~J.~Salzer, and E.~D.~Skillman,
Improving helium abundance determinations with Leo P as a case study,
JCAP \textbf{03}, 027 (2021);
arXiv:2010.04180.

\bibitem{EMPRESS2022}
A.~Matsumoto \emph{et~al.} (EMPRESS Collaboration),
EMPRESS. VIII. A new determination of primordial helium abundance with extremely metal-poor
galaxies: a suggestion of the lepton asymmetry and implications for the Hubble tension,
Astrophys.\ J.\ \textbf{941}, 167 (2022);
arXiv:2203.09617.

\end{thebibliography}

\end{document}
